--- docs/design/proposal/manuscript/sections/corporate-methods.tex

+++ docs/design/proposal/manuscript-intragroup-fx-20261013/sections/corporate-methods.tex

@@ -66,4 +66,6 @@

 The retained AASB 116.7/15, 29--30 and 48--62 support cost and depreciation. Asset 100, residual 10, three-year even use gives 30 per full year. A year-two estimate change after first-year depreciation to residual 10 and three remaining years gives (70\ensuremath{-}10)/3=20 prospectively; an error in original cost is a correction, not an estimate change. Component replacements derecognize the replaced component's remaining amount before recognizing its replacement; no perpetual double asset. Borrowing-cost capitalization, impairment, provisions, employee benefits and share-based payment are specified as additional scoped methods below, not silently assigned zero in real-company cases.
 
+\input{sections/apple-operating-assets}
+
 \hypertarget{lease-and-tax-integration}{%
 \section{Lease and tax integration}\label{method-corporate-methods-lease-and-tax-integration}}
@@ -93,4 +95,6 @@

 2024 assets 811.5 equal liabilities 181 plus equity 630.5; 2025 assets 813.25 equal liabilities 44.75 plus equity 768.5. For this case, elect financing classification for lease interest under the retained AASB 107 regime. The operating and financing cash totals below depend on that presentation election. Direct operating cash is 300 then 255.25 after the first tax payment; financing lease cash is \ensuremath{-}121 each year; no investing cash arises. Indirect operating cash in year one is profit 130.5 plus depreciation 105 plus interest 21 minus deferred benefit 1.25 plus unpaid current tax 44.75 = 300. Year two gives profit 138 +105 +11 +1.25 =255.25. The second tax payable settles next year. Zero OCI and no owner transactions mean the equity change equals cumulative profit. This case tests gross tax, cash timing and financial-statement coherence, not a real entity's tax law.
 
+\input{sections/apple-tax-settlement}
+
 \hypertarget{returning-to-the-company-forecast}{%
 \section{Returning to the company forecast}\label{method-corporate-methods-returning-to-the-company-forecast}}
